% Generated draft: requires mathematical and source review.
\documentclass[11pt]{article}
\usepackage[margin=25mm]{geometry}
\usepackage{fontspec}
\setmainfont{Latin Modern Roman}
\usepackage{amsmath,amssymb,mathtools,unicode-math}
\setmathfont{Latin Modern Math}
\usepackage{enumitem,needspace,xurl,hyperref}
\setlength{\parindent}{0pt}
\setlength{\parskip}{5pt}
\setlength{\emergencystretch}{3em}
\newcommand{\N}{\mathbb{N}}
\newcommand{\Z}{\mathbb{Z}}
\newcommand{\Q}{\mathbb{Q}}
\newcommand{\R}{\mathbb{R}}
\newcommand{\C}{\mathbb{C}}
\newcommand{\F}{\mathbb{F}}
\newcommand{\A}{\mathbb{A}}
\newcommand{\PP}{\mathbb{P}}
\newcommand{\E}{\mathbb{E}}
\newcommand{\eps}{\varepsilon}
\newcommand{\dd}{\,\mathrm{d}}
\begin{document}
\section*{UnsolvedMath ID: 700782}
\textbf{Source problem code:} catalogue entry
\subsection{Definitions and mathematical statement}
Let \(\mathfrak V\) be a variety of groups, that is, a class of groups defined by group identities, and let \(k\) be an infinite cardinal. For this draft, interpret rank as generating rank:
\[
\operatorname{rk}(G)=\min\{\lvert X\rvert:X\subseteq G,\ \langle X\rangle=G\}.
\]
The supplied source does not explicitly define rank, so this convention requires confirmation. A group \(F\in\mathfrak V\) is free in \(\mathfrak V\) if it has a generating set \(X\) such that, for every \(L\in\mathfrak V\), each map \(X\to L\) extends uniquely to a homomorphism \(F\to L\).

As defined in the source, a group \(G\in\mathfrak V\) with \(\operatorname{rk}(G)=k\) is almost free in \(\mathfrak V\) if
\[
\forall H\leq G,\quad
\operatorname{rk}(H)<k\ \Longrightarrow\
\exists K\leq G\quad
\bigl(H\leq K\ \text{and }K\text{ is free in }\mathfrak V\bigr).
\]
There is no stated requirement that \(\operatorname{rk}(K)<k\), or that \(K\) be a free factor of \(G\).

The source asks for which infinite cardinals \(k\) there exist almost free groups of rank \(k\) that are not free, but its final sentence uses \(\mathfrak B\) where its definition uses \(\mathfrak V\). No meaning or relationship is supplied for \(\mathfrak B\). Consequently, the literal target is unresolved. If \(\mathfrak B\) was intended to denote \(\mathfrak V\), the question is: for a given variety \(\mathfrak V\), determine the infinite cardinals \(k\) for which
\[
\exists G\in\mathfrak V\quad
\left[
\begin{gathered}
\operatorname{rk}(G)=k,\qquad G\text{ is not free in }\mathfrak V,\\
\forall H\leq G,\quad
\operatorname{rk}(H)<k\ \Longrightarrow\
\exists K\leq G\quad
\bigl(H\leq K\ \text{and }K\text{ is free in }\mathfrak V\bigr)
\end{gathered}
\right].
\]
This displayed question is conditional on identifying the two variety symbols; that identification is not established by the supplied record.
\subsection{Short English statement}
For a given variety of groups, determine which infinite ranks admit a group that is not relatively free but whose every subgroup of smaller rank lies in a relatively free subgroup. The source changes its variety symbol from \(\mathfrak V\) to an undefined \(\mathfrak B\), so the intended target needs confirmation.
\subsection{Sources}
\begin{itemize}
\item Local UnsolvedMath registry, supplied source record: permanent ID \(700782\), exact source problem code \texttt{catalogue entry}; upstream revision \texttt{d371ee917ac9647fbb06cd6fc529deec6f3dda82}.
\item The supplied attribution is to S.\ V.\ Rychkov, proposed in 1990; \emph{The Kourovka Notebook (Unsolved Problems in Group Theory)}, original problem 11.90, 11th issue (1990). The record identifies the 21st issue edition (2026), \texttt{arXiv:1401.0300v48}, and names E.\ I.\ Khukhro and V.\ D.\ Mazurov as editors of the collection.
\item The registry labels this problem \emph{open}. This is a supplied status, not an independently verified mathematical or literature assessment. The cited source has not been independently consulted for this draft.
\end{itemize}
\end{document}
